%%% FIELD relative-proof
Fix an arm \(d\), and suppress the arm subscript when no ambiguity can
result.  Write \(J=|J_d|\), \(M=M_d\), \(A=A_d\), \(m=m_d\),
\(U=U_d\), and \(T(g)=\int_B g\,d\nu\).  All norms on the moment
space below are Euclidean.  The proof has four parts: a geometric projection
argument, a compactness argument, a support-function moment correction, and a
clipping-preserving finite-element approximation.

First consider the projections.  Because the initial triangulation conforms to
every cell boundary and to \(B\), so does every midpoint refinement.  If
\(x\in C_{dj}\), some triangle lying in \(C_{dj}\) and containing \(x\) has a
vertex in \(\mathcal N_h\cap C_{dj}\) at distance at most its diameter.  The
same argument on a boundary edge applies to \(x\in B\).  Nearest-node
selection, with lexicographic tie breaking, is Borel measurable and therefore
satisfies
\[
 \|x-p_{dj,h}(x)\|_2\le h,
 \qquad
 \|x-p_{B,h}(x)\|_2\le h.                                      \tag{1}
\]
Here and below the support condition on \(\mu_{dj}\) is used when integrating
the first inequality.  For any \(g\in\mathcal F_d\), set
\(v_z=g(z)\) at every \(z\in\mathcal N_h\).  The range restriction and the
\(M\)-Lipschitz restriction give all nodal constraints.  The pushforward
definitions and (1) give, for every \(j\in J_d\),
\begin{align}
 \left|\sum_z v_z\mu_{dj,h}(\{z\})-m_{dj}\right|
 &=\left|\int_{C_{dj}}\{g(p_{dj,h}(x))-g(x)\}\,d\mu_{dj}(x)\right|
 \le Mh,                                                       \tag{2}\\
 \left|\sum_z v_z\nu_h(\{z\})-T(g)\right|
 &=\left|\int_B\{g(p_{B,h}(x))-g(x)\}\,d\nu(x)\right|
 \le Mh.                                                       \tag{3}
\end{align}
Both laws in these displays have total mass one.  Compatibility makes the
exact moment fiber nonempty.  It is uniformly closed inside the uniformly
bounded, common-modulus class \(\mathcal F_d\), and that class is compact in
the uniform norm: to see this without importing a compactness theorem, choose
successive finite \(1/k\)-nets of compact \(S\), diagonalize a bounded sequence
on the countable union of the nets, and use the common Lipschitz modulus to
upgrade pointwise convergence on that dense set to uniform convergence.  The
uniform limit remains in \(\mathcal F_d\).  Thus an optimizer \(g^*\) attaining
\(U(m)\) exists.  Applying (2)--(3) to \(g^*\) proves
\[
 U(m)\le V_{d,h}(m)+Mh.
\]
Since \(0\le U(m)\le1\), the clipping in the definition of \(B_{d,h}\)
preserves the inequality:
\[
 U(m)\le B_{d,h}(m).                                         \tag{4}
\]

Conversely, let \(v\) be feasible for the atomic program.  Pairwise nodal
feasibility implies
\[
 \min_{z'\in\mathcal N_h}\{v_{z'}+M\|z-z'\|_2\}=v_z
 \quad(z\in\mathcal N_h).
\]
Indeed the left side is at most \(v_z\) by choosing \(z'=z\), and every other
term is at least \(v_z\) by the pairwise constraint.  The minimum of
\(M\)-Lipschitz functions is \(M\)-Lipschitz, and clipping by
\(t\mapsto\min\{1,\max\{0,t\}\}\) is a contraction.  Hence
\(e_h=E_{d,h}v\) belongs to \(\mathcal F_d\) and equals \(v\) at every node.
Combining the relaxed moment constraint with (1) yields
\[
 \|Ae_h-m\|_2\le 2Mh\sqrt J,                                 \tag{5}
\]
while the boundary calculation yields
\[
 \left|T(e_h)-\sum_zv_z\nu_h(\{z\})\right|\le Mh.            \tag{6}
\]
The finite atomic feasible set is a nonempty compact polytope, so an atomic
optimizer exists.  Along an arbitrary sequence \(h_n\downarrow0\), take such
optimizers and their extensions.  The elementary compactness argument above
gives a uniformly convergent subsequence with limit \(e\in\mathcal F_d\).
Equations (5)--(6) imply \(Ae=m\) and that the corresponding atomic objectives
converge to \(T(e)\le U(m)\).  Since every subsequence has a further subsequence
with this upper-limit property,
\(\limsup_{h\downarrow0}V_{d,h}(m)\le U(m)\).  Moreover
\(0\le B_{d,h}(m)-V_{d,h}(m)\le Mh\).  Together with (4), this proves
\[
 B_{d,h}(m)\longrightarrow U(m)                              \tag{7}
\]
at every compatible mean, including a mean on the relative boundary of the
moment body.  This step is a value-side compactness argument and makes no exact
finite-mesh feasibility claim.

We next derive the relative-interior quantitative bound.  Put
\(r=r_d(m)>0\) and \(q=Ae_h\).  Since \(q,m\in K_d\), the vector
\(q-m\) belongs to \(E_d\).  If \(q\ne m\), put
\(a=(q-m)/\|q-m\|_2\).  The definition of \(r\) gives
\(m-ra\in K_d\), and the identity
\[
 m=\frac{r}{r+\|q-m\|_2}q
   +\frac{\|q-m\|_2}{r+\|q-m\|_2}(m-ra)                     \tag{8}
\]
holds.  The endpoint map \(U\) is concave on \(K_d\): convex combinations of
two feasible surfaces have the corresponding convex-combination moments and
target values.  Applying concavity to (8) and using \(0\le U\le1\) gives
\[
 U(q)\le U(m)+\frac{\|q-m\|_2}{r}.                           \tag{9}
\]
The same inequality is immediate if \(q=m\).  For an atomic optimizer,
(6), feasibility of \(e_h\) at moment \(q\), and (9) imply
\[
 V_{d,h}(m)\le T(e_h)+Mh
 \le U(m)+\frac{\|q-m\|_2}{r}+Mh.
\]
Adding the final \(Mh\) in the definition of \(B_{d,h}\), noting that clipping
at one can only decrease the result, and using (5), proves the stated exact
outer error:
\[
 0\le B_{d,h}(m)-U(m)
 \le 2Mh\left(1+\frac{\sqrt J}{r}\right).                   \tag{10}
\]

For the inner correction, set
\[
 K_{d,h}=A(\mathcal F_{d,h}),
 \qquad
 \delta_h=\sqrt J\,\varepsilon_{d,h}.
\]
The nodal feasible set defining \(\mathcal F_{d,h}\) is compact and convex.
Vertex range bounds extend to each triangle by barycentricity.  On any line
segment in convex \(S\), integration of the piecewise-constant gradient and
the trianglewise gradient bound give the global Euclidean Lipschitz inequality.
Thus
\[
 \mathcal F_{d,h}\subseteq\mathcal F_d,
 \qquad K_{d,h}\subseteq K_d,                               \tag{11}
\]
and \(K_{d,h}\) is compact and convex.  Let \(s_C(a)=\sup_{x\in C}a^\top x\)
denote a support function.  For \(a\in E_d\), choose a support maximizer in
\(\mathcal F_d\) and a best uniform approximant in \(\mathcal F_{d,h}\).
Because every \(\mu_{dj}\) is a probability law,
\(\|A(g-q)\|_2\le\sqrt J\|g-q\|_\infty\).  Therefore
\begin{align}
 s_{K_{d,h}}(a)
 &\ge s_{K_d}(a)-\delta_h\|a\|_2\\
 &\ge a^\top m+(r-\delta_h)\|a\|_2.                        \tag{12}
\end{align}
The second inequality uses
\(m+r\mathbb B_{E_d}\subseteq K_d\).  Equation (12) implies
\[
 m+(r-\delta_h)\mathbb B_{E_d}\subseteq K_{d,h}             \tag{13}
\]
whenever \(\delta_h<r\).  For completeness, if a point of the ball in (13)
were outside the compact convex set \(K_{d,h}\), its nearest point in the
finite-dimensional affine space \(m+E_d\) would furnish a separating vector
\(a\in E_d\), contradicting (12).  This is the support-function separation
step.

Let \(g^*\) attain \(U(m)\), and choose
\(g_h\in\mathcal F_{d,h}\) with
\(\|g_h-g^*\|_\infty\le\varepsilon_{d,h}\).  Put
\(e=Ag_h-m\); then \(e\in E_d\) and \(\|e\|_2\le\delta_h\).  If
\(\delta_h=0\), the surface \(g_h\) is exactly moment feasible.  Otherwise set
\(t=\delta_h/r\).  Since
\[
 \left\|-\frac{1-t}{t}e\right\|_2
 \le r-\delta_h,
\]
(13) supplies \(q_h\in\mathcal F_{d,h}\) such that
\[
 Aq_h=m-\frac{1-t}{t}e.
\]
The convex combination \(f_h=(1-t)g_h+tq_h\) then satisfies
\(Af_h=m\).  Positivity and normalization of \(\nu\), together with the range
restriction, imply \(0\le T(q_h),T(g^*)\le1\).  Hence
\begin{align}
 T(f_h)
 &=(1-t)T(g_h)+tT(q_h)\\
 &\ge(1-t)\{U(m)-\varepsilon_{d,h}\}\\
 &\ge U(m)-\varepsilon_{d,h}-\frac{\delta_h}{r}.            \tag{14}
\end{align}
Since \(\delta_h<r\) for all sufficiently small \(h\), this proves eventual
nonemptiness of the exact piecewise-affine fiber.  Inclusion (11) gives
\(P_{d,h}(m)\le U(m)\), while (14) gives
\[
 0\le U(m)-P_{d,h}(m)
 \le\varepsilon_{d,h}\left(1+\frac{\sqrt J}{r}\right).     \tag{15}
\]

It remains to prove, rather than assume, the claimed approximation rate.  All
children of a triangle in a midpoint refinement are similar to their parent.
Consequently, for the edge-row matrix \(R_T\) of a triangle \(T\),
\[
 C_0:=\sup_{h,T\in\mathcal T_h}
       \frac{\operatorname{diam}(T)\|R_T^{-1}\|_{\mathrm{op}}}{\sqrt2}<\infty.
                                                                    \tag{16}
\]
Suppose first that \(M>0\).  The clipped McShane formula extends
\(g\in\mathcal F_d\) to an \([0,1]\)-valued \(M\)-Lipschitz function
\(\widetilde g\) on \(\mathbb R^2\).  For a standard Gaussian vector
\(Z\in\mathbb R^2\), define
\[
 u_\delta(x)=\mathbb E\{\widetilde g(x+\delta Z)\}.
\]
Direct coupling and the polar density
\(s e^{-s^2/2}\mathbf 1\{s\ge0\}\) give
\[
 \|u_\delta-g\|_\infty
 \le M\delta\,\mathbb E\|Z\|_2
 =M\delta\sqrt{\frac\pi2}.                                \tag{17}
\]
The function \(\widetilde g\) has distributional first derivatives bounded in
Euclidean norm by \(M\): this follows directly by taking bounded directional
difference quotients, testing against a smooth compactly supported function,
and passing to a weak limit in the integration-by-parts identity.  Convolution
with the Gaussian kernel is smooth.  Differentiating the convolution of that
bounded weak gradient in a unit direction \(a\) gives
\[
 D_a\nabla u_\delta(x)
 =\frac1\delta\mathbb E\!\left[
       \nabla\widetilde g(x+\delta Z)(Z^\top a)\right].
\]
Since \(\mathbb E|Z^\top a|=\sqrt{2/\pi}\),
\[
 \|D^2u_\delta\|_{\mathrm{op}}
 \le\frac{M}{\delta}\sqrt{\frac2\pi}.                     \tag{18}
\]

Let \(I_hu_\delta\) be the nodal affine interpolant.  Taylor expansion from one
vertex along the other two edges shows that the two-component remainder vector
has norm at most
\(\|D^2u_\delta\|_{\mathrm{op}}\operatorname{diam}(T)^2/\sqrt2\).
Multiplication by \(R_T^{-1}\), followed by (16)--(18), yields
\[
 \|\nabla(I_hu_\delta)|_T\|_2\le M(1+\vartheta),
 \qquad
 \vartheta=C_0\sqrt{\frac2\pi}\frac h\delta.              \tag{19}
\]
Barycentric interpolation preserves the range and, using only that
\(u_\delta\) is \(M\)-Lipschitz, also gives
\[
 \|I_hu_\delta-u_\delta\|_\infty\le Mh.                    \tag{20}
\]
Contract the interpolant about \(1/2\):
\[
 q_h=\frac12+\frac{I_hu_\delta-1/2}{1+\vartheta}.
\]
Equations (19)--(20) show that \(q_h\in\mathcal F_{d,h}\).  Since
\(|I_hu_\delta-1/2|\le1/2\), equations (17)--(20) imply
\[
 \|q_h-g\|_\infty
 \le M\delta\sqrt{\frac\pi2}+Mh+\frac{\vartheta}{2}.       \tag{21}
\]
The right side is uniform over \(g\in\mathcal F_d\).  Taking
\(\delta=\sqrt h\) in (21) proves
\[
 \varepsilon_{d,h}=O(\sqrt h),                              \tag{22}
\]
where the constant depends only on \(M_d\) and the fixed initial mesh.  If
\(M=0\), convexity of \(S\) makes every class member constant, and constants
belong to every \(\mathcal F_{d,h}\); hence
\(\varepsilon_{d,h}=0\).  This completes the inner-program proof.

Finally, \(g\mapsto1-g\) bijects the fiber at \(m_d\) with the fiber at
\(\mathbf1-m_d\), preserves every approximation class, and gives
\(L_d(m_d)=1-U_d(\mathbf1-m_d)\).  Applying (4), (7), (10), and (15) to the
reflected programs supplies the lower-endpoint bounds.  Combining the two
armwise endpoint brackets by interval subtraction gives an outer bracket and
an inner attainable interval for \(I(m_0,m_1)\); their endpoint gaps are sums
of the corresponding armwise gaps and therefore vanish when the certified
finite-program optimization errors vanish.  This proves every assertion of the
theorem.

%%% FIELD terminating-proof
Fix an arm and again suppress its subscript.  We construct a terminating
procedure; no efficiency or bit-complexity assertion is needed.  The technique
is dovetailed enumeration of exact rational witnesses and exact algebraic
dual certificates.  Termination follows from the relative-interior program
theorem and finite multiplier attainment in the clipping-aware dual theorem.

We first describe the exact data representation.  On a rational conforming
triangulation, a continuous piecewise-affine function is represented by its
nodal vector \(z\).  Conformity to all density breaks implies
\[
 Az=C_hz,                                                     \tag{23}
\]
where \(C_h\) is a rational matrix: on each rational triangle the integral of
each barycentric coordinate is one third of the rational triangle area, and
the density is rational and constant.  Nodal clipping is a finite list of
rational linear inequalities.  If \(R_T\) is the rational edge matrix of a
triangle, then \(\nabla(z|_T)=R_T^{-1}\Delta_Tz\) has rational coordinates
when \(z\) is rational, and
\[
 \|\nabla(z|_T)\|_2^2\le M^2                               \tag{24}
\]
is an exactly decidable rational inequality.  Exact equality in (23), rather
than a floating residual, is part of every accepted primal certificate.

On a rational boundary subedge of length \(\ell\), with endpoint function
values \(z_0,z_1\) and endpoint weight values \(w_0,w_1\), direct integration
of the product of two affine functions gives
\[
 \int_e zw\,d\mathcal H^1
 =\ell\,
   \frac{(2w_0+w_1)z_0+(w_0+2w_1)z_1}{6}.                  \tag{25}
\]
Thus a rational nodal vector has an exactly represented real-algebraic boundary
objective.  Rational intervals enclosing (25) to any prescribed width are
obtained by integer-arithmetic bisection of the finitely many square roots.

We now prove density of exact rational primal certificates.  Put
\(c=\mathbf1/2=A(1/2)\).  Suppose \(M>0\).  If \(m\ne c\), choose a sufficiently
small rational \(s\in(0,1)\) and set
\[
 m'=\frac{m-sc}{1-s}
    =m+\frac{s}{1-s}(m-c).                                  \tag{26}
\]
Because \(m\in\operatorname{ri}K_d\), the relative ball in the definition of
\(r_d(m)\) contains (26) for all sufficiently small \(s\), and then
\(m'\in\operatorname{ri}K_d\).  The relative-interior program theorem supplies,
on a sufficiently fine mesh, \(q'\in\mathcal F_{d,h}\) with \(Aq'=m'\).
Consequently
\[
 q^0=(1-s)q'+s/2,
 \qquad Aq^0=m,                                              \tag{27}
\]
has every nodal value in \([s/2,1-s/2]\) and every triangle gradient norm at
most \((1-s)M<M\).  If \(m=c\), the constant \(q^0=1/2\) has the same strict
properties.  Hence the exact moment fiber contains a point strict in all range
and gradient inequalities.

Given \(\eta>0\), the same theorem supplies an exactly feasible
piecewise-affine \(p\) with \(T(p)>U(m)-\eta/4\).  Passing both \(p\) and
\(q^0\) to a common nested refinement and mixing an arbitrarily small positive
amount of \(q^0\) into \(p\) produces a strict point \(p^0\) with
\[
 Ap^0=m,
 \qquad T(p^0)>U(m)-\eta/2.                                 \tag{28}
\]
The affine solution space \(\{z:C_hz=m\}\) is defined over \(\mathbb Q\),
because \(C_h\) and the assumed input \(m\) are rational.  Rational Gaussian
elimination gives a rational particular solution and a rational basis of its
nullspace.  Approximating the free coordinates of \(p^0\) by rationals gives
rational points in the same exact affine space converging to \(p^0\).  The
strict inequalities persist, and (25) is continuous.  Therefore the
enumeration eventually finds a rational nodal \(q^U\) such that
\[
 q^U\in\mathcal F_d,
 \qquad Aq^U=m,
 \qquad T(q^U)>U(m)-\eta.                                   \tag{29}
\]
Every assertion in (29) except the last comparison with the unknown endpoint
is certified by (23)--(25); the comparison is used only to prove termination.
If \(M=0\), every feasible surface is constant.  Compatibility makes
\(m=c_0\mathbf1\); rationality of \(m\) makes \(c_0\) rational, and the constant
surface gives both endpoints exactly.  This handles the zero-radius case.

Next construct the dual enclosure for \(M>0\).  By the clipping-aware dual
theorem and \(r_d(m)>0\), there is a finite optimal multiplier
\(\lambda^*\in\mathbb R^J\).  The objective is continuous.  More explicitly,
for any \(\lambda,\lambda'\), the range restriction and the fact that
\(m_j\in[0,1]\) imply
\begin{align}
 |D(\lambda;m)-D(\lambda';m)|
 &\le |(\lambda-\lambda')^\top m|
   +\sup_{g\in\mathcal F_d}
       \left|\sum_j(\lambda_j-\lambda'_j)\int g\,d\mu_{dj}\right|\\
 &\le2\|\lambda-\lambda'\|_1.                              \tag{30}
\end{align}
Thus enumeration of rational multipliers contains multipliers with dual value
arbitrarily close to \(U(m)\).

For a rational multiplier \(\lambda\), form
\[
 \sigma(\lambda)=\nu-\sum_j\lambda_j\mu_{dj},
 \qquad
 \sigma_h(\lambda)=\nu_h-\sum_j\lambda_j\mu_{dj,h}.        \tag{31}
\]
The nearest-node Voronoi boundaries are rational lines, because the difference
of squared distances to two rational nodes is rational affine.  Their
intersections with rational triangles and rational boundary segments therefore
have rational vertices.  Cell atomic masses are rational.  Boundary atomic
masses, Euclidean node distances, and the quantities in (25) belong to a finite
multiquadratic extension
\[
 \mathbb Q(\sqrt{s_1},\ldots,\sqrt{s_R})                    \tag{32}
\]
for squarefree positive integers \(s_r\).  They are not silently treated as
rational.  Exact addition and multiplication use the finite basis of products
of independent square roots.  Equality is decided by reducing in that basis;
the sign of a nonzero element is decided by rational isolating intervals,
whose separation terminates because the element is algebraic and nonzero.
This supplies terminating exact sign tests for every atom of \(\sigma_h\), and
hence an exact Jordan decomposition.

Enumerate finite matrices \(\pi=(\pi_{xy})\) with nonnegative rational entries
and use the exact comparisons just described to check
\[
 \sum_y\pi_{xy}\le\sigma_h^+(\{x\}),
 \qquad
 \sum_x\pi_{xy}\le\sigma_h^-(\{y\}).                       \tag{33}
\]
For every matrix passing (33), the attained partial-transport identity in the
clipping-aware dual theorem gives
\[
 H_M(\sigma_h)
 \le \sigma_h^+(S)-\sum_{x,y}\pi_{xy}
      +M\sum_{x,y}\|x-y\|_2\pi_{xy}.                        \tag{34}
\]
The right side is algebraic and has computable rational outward enclosures.
These rational plans approximate the atomic optimum: take an attained optimal
real plan and approximate every positive entry from below by rationals.  The
result remains feasible in (33), converges entrywise to the optimal plan, and
the finite objective in (34) converges to the optimum.  This argument also
covers zero capacities without waiting for an inconclusive floating sign.

The continuous remainder is an analytic inequality, not a mesh residual.  For
every \(g\in\mathcal F_d\), equations (1) and (31) give
\begin{align}
 &\left|\int g\,d\sigma(\lambda)-\int g\,d\sigma_h(\lambda)\right|\\
 &\quad\le
 \left|\int_B\{g(x)-g(p_{B,h}(x))\}\,d\nu(x)\right|
 +\sum_j|\lambda_j|
   \left|\int\{g(x)-g(p_{dj,h}(x))\}\,d\mu_{dj}(x)\right|\\
 &\quad\le Mh(1+\|\lambda\|_1).                            \tag{35}
\end{align}
Consequently, if \(\mathcal C_h(\lambda,\pi)\) denotes the right side of (34), any
rational upward rounding of
\[
 \lambda^\top m+\mathcal C_h(\lambda,\pi)+Mh(1+\|\lambda\|_1)\tag{36}
\]
is a certified upper bound for \(U(m)\).  A rational downward rounding of
\(T(q^U)\) is a certified lower bound by (29).  The maximum mesh diameter and
all square roots in (36) are themselves algebraic and are outwardly rounded by
the same exact isolation procedure.

Define \(\mathsf{VPA}\) to dovetail over the midpoint-refinement level; bounds
on numerators and denominators of rational nodal vectors, multipliers, and plan
entries; and the precision of all rational isolating intervals.  At each finite
stage it checks (23), the nodal range inequalities, (24), (33), and the rational
ordering of the proposed lower and upper bounds.  Every such check terminates.
It halts when their gap is at most the requested rational \(\epsilon\).

To prove that it halts, fix \(\eta>0\).  Equation (29) supplies a rational
primal candidate with lower objective greater than \(U(m)-\eta\).  Equation
(30) supplies a rational \(\lambda\) with
\(D(\lambda;m)<U(m)+\eta\).  Keep this multiplier fixed.  By (35), sufficiently
fine refinement makes its remainder smaller than \(\eta\).  For that finite
atomic problem, rational-plan density in (34) supplies a cost within \(\eta\)
of its optimum.  Since (35) also bounds the atomic support value above by the
continuous support value plus \(\eta\), expression (36) is then less than
\(U(m)+4\eta\).  Choose all upward and downward rational rounding errors below
\(\eta\).  The resulting certified gap is less than \(7\eta\).  Taking
\(7\eta<\epsilon\), the dovetailed enumeration must eventually encounter this
finite collection of candidates and halt.  Notice that neither \(r_d(m)\) nor
an unknown convergence rate has to be supplied to the algorithm; the
relative-interior condition is used to prove that a terminating candidate
exists.

Reflection completes the lower endpoint.  The map \(g\mapsto1-g\) sends
\(K_d\) onto \(\mathbf1-K_d\), sends \(m\) to \(\mathbf1-m\), and preserves the
relative-interior radius.  Run the same procedure at \(\mathbf1-m\).  If
\([\underline U',\overline U']\) encloses
\(U_d(\mathbf1-m)\), then
\[
 [1-\overline U',1-\underline U']
\]
encloses \(L_d(m)=1-U_d(\mathbf1-m)\).  If \(q'\) is the exact upper primal
witness in the reflected problem, set \(q^L=1-q'\); then
\(Aq^L=m\).  This returns the two claimed rational-nodal surfaces and both
rational outward endpoint enclosures with gaps at most \(\epsilon\).  Equations
(23)--(25), (33)--(36) are precisely the advertised exact certificate checks.

We finally give the requested rational two-cell strict-interior witness.  Let
\[
 S=[-1,1]\times[0,1],\quad B=\{0\}\times[0,1],
\]
let
\(C_1=[0,1]\times[0,1/2]\) and
\(C_2=[0,1]\times[1/2,1]\), give each cell its uniform law, let \(\nu\) be
uniform on \(B\), and take \(M=1/2\) and
\(m=(1/2,1/2)\).  Rational piecewise-affine hat functions supported strictly
inside the respective cells have a diagonal, strictly positive two-by-two
moment matrix.  Sufficiently small signed combinations of those hats with the
constant \(1/2\) retain strict clipping and Lipschitz slack.  Hence a Euclidean
neighborhood of \(m\) lies in \(K_d\), so \(r_d(m)>0\).

The mixture \(\rho=(\mu_1+\mu_2)/2\) is uniform on the unit square.  Coupling
\((x,y)\) to \((0,y)\) and using the Lipschitz inequality gives, for every
feasible \(g\),
\[
 T(g)-\frac12\sum_{j=1}^2\int g\,d\mu_j
 \le M\int_0^1\int_0^1x\,dx\,dy=\frac14.                   \tag{37}
\]
The globally feasible rational piecewise-affine functions
\[
 q^U(x,y)=\frac34-\frac12\max\{x,0\},
 \qquad
 q^L(x,y)=\frac14+\frac12\max\{x,0\}                       \tag{38}
\]
have both cell means exactly \(1/2\), and their boundary means are respectively
\(3/4\) and \(1/4\).  Thus (37) and reflection prove
\[
 U(m)=\frac34,
 \qquad L(m)=\frac14.                                      \tag{39}
\]
The exact multiplier \(\lambda=(1/2,1/2)\) has
\(\lambda^\top m=1/2\).  On refinement level \(k\), let
\(\bar h_k\in\mathbb Q\) be an outward upper bound on the maximum triangle
diameter, with \(\bar h_k\downarrow0\).  Pushing the coupling in (37) through
the cell and boundary projections increases each transported distance by at
most \(2\bar h_k\).  The induced finite transportation problem has rational
capacities in this axis-aligned dyadic example; canceling common atomic mass
and taking a rational basic flow gives a partial-transport certificate with
\[
 H_M\!\left(\nu_h-\frac12\mu_{1,h}-\frac12\mu_{2,h}\right)
 \le\frac14+\bar h_k.                                      \tag{40}
\]
The continuous remainder in (35) is
\(M\bar h_k(1+\|\lambda\|_1)=\bar h_k\).  Equations
(38)--(40) therefore give the exactly certified rational enclosures
\[
 \frac34\le U(m)\le\frac34+2\bar h_k,
 \qquad
 \frac14-2\bar h_k\le L(m)\le\frac14.                      \tag{41}
\]
Their gaps shrink to zero, the primal moments in (38) are exact continuous
cell integrals, and no floating residual is used.  This passes the stipulated
early-kill test and completes the construction.

%%% FIELD oeq-answer-statement
Let \(S\) be a compact convex rational polygon, let \(B\) and the cells be
rational polygonal, let every \(\mu_{dj}\) have a rational piecewise-constant
density, and let \(d\nu=w\,d\mathcal H^1\), where \(w\) is nonnegative,
normalized, and rational piecewise affine on \(B\).  For rational \(M_d\ge0\),
rational \(m_d\in K_d\) with \(r_d(m_d)>0\), and rational \(\epsilon>0\), the
answer is affirmative.  The procedure \(\mathsf{VPA}\) terminates with
rational-nodal continuous piecewise-affine witnesses in \(\mathcal F_d\) whose
\(A_d\)-moments equal \(m_d\), and with rational outward enclosures of
\(U_d(m_d)\) and \(L_d(m_d)\), each of width at most \(\epsilon\).  Its
certificate checks exact continuous cell integrals, clipping, squared Euclidean
gradient inequalities, a feasible atomic partial-transport plan, and the
continuous remainder \(M_dh(1+\|\lambda_d\|_1)\); no floating-residual or
bit-complexity claim is made.

%%% FIELD oeq-answer-proof
Apply the terminating verified-reconstruction theorem to the stated rational
input.  Its hypotheses are exactly the geometry, density, normalization,
compatibility, positive relative-interior radius, and finite rational tolerance
listed here.  That theorem constructs the required two primal witnesses and
the two rational endpoint enclosures, proves termination by dovetailed exact
enumeration, and verifies each listed primal, transport, and continuous-remainder
condition.  Therefore it gives the affirmative answer verbatim.
